Each cell is minimised in two stages, to a force norm of $10^{-3}$ and then
$3\times10^{-4}$~eV/\AA. Between the stages the stresses in shells more than
5~\AA{} from the inclusion surface change by at most «G1_MAX»~MPa.

The $+\varepsilon$ cell starts the minimisation from the elastic solution and
the $-\varepsilon$ cell from the atom positions of the reference cell. Half
the sum of the responses to $+\varepsilon$ and $-\varepsilon$ therefore
measures both nonlinearity and the dependence on the starting point. In
shells more than 5~\AA{} from the surface it is at most «G3_MAX»~MPa, about
1\% of the stresses along and across the field. The main text uses half
the difference of the two responses.

No matrix atom jumps to a new position. The number of matrix atoms whose
displacement differs from the elastic solution by more than 0.1~\AA{} is
«EV_PLUS» in the $+\varepsilon$ cell and «EV_MINUS» in the $-\varepsilon$ cell.
All such atoms lie within «EV_MINUS_D»~\AA{} of the surface and depart by at
most «EV_MINUS_MAX»~\AA; no atom departs this much in both cells.

Defects are searched for with DXA (the dislocation extraction algorithm)
\cite{Stukowski2012DXA}, which finds dislocation lines from the distortion of
the lattice around them, and with a nearest-neighbour analysis, which
identifies the crystal surroundings of each atom. The search covers the
reference and $\pm\varepsilon$ cells, after each minimisation stage, with and
without the inclusion atoms. No dislocation segment is found in any cell. The
aluminium atoms without FCC (face-centred cubic) surroundings more than 6~\AA{}
from the surface are nearly the same in all three cells: «NONFCC_FAR» atoms
6--7.5~\AA{} from it, at the edge of the interface layer.

The mean added stress over the matrix atoms ($xx$, $yy$, $zz$) is
$«G4_MD»$~MPa; the elastic solution for the same cell gives $«G4_CT»$~MPa.
Their ratio, about «G4_RATIO», is close to the factor «A_FIT» with which the
atomistic field matches the elastic solution (Section~3.2 of the main text).

The atomistic field is weaker than the elastic solution because of a compliant
aluminium layer at the interface. The displacements
of the matrix atoms, relative to those of the elastic solution for the same
cell, are «KIN_RAD»\% along the surface normal and «KIN_TAN»\% along the
surface in the first layer ($d<2$~\AA), «D46R»\% and «D46T»\% at
$d=4$--6~\AA, and «D1214R»\% and «D1214T»\% at $d=12$--14~\AA. The strain of
each atom is obtained from the displacements of its neighbours within 3.2~\AA{} (the twelve nearest in FCC).
For $d>8$~\AA{} the virial stress equals the cubic elastic constants of the
potential times this strain with a slope of «VIR_FAR» (correlation at least
«VIR_CORR»); for $d<3$~\AA{} the slope is «VIR_NEAR», so the layer at the
interface carries less stress for the same strain. The aluminium density at
$r=25.0$--$25.5$~\AA{} is «DENS_FIRST» of bulk. The cavity left for the
inclusion after removing the aluminium atoms within 2.2~\AA{} of it is
«CAV_PCT»\% larger than the inclusion in volume. The far field alone does not
tell such a compliant layer from a perfectly bonded sphere of smaller radius:
the factor «A_FIT» corresponds to a radius of «R_EQ»~\AA. The tangential
displacement in the first layer is «KIN_TAN»\% of the elastic one, so at 0~K
the interface does not slide freely.

With the inclusion strain 4 and 10 times the given one (0.78 and 1.94\%,
minimised to $10^{-3}$~eV/\AA) there are no defects either: DXA finds no
dislocations, there are no more atoms without FCC surroundings beyond 6~\AA{}
from the surface than in the reference cell, and the field divided by the
strain multiple equals the field at $\varepsilon$ (median ratios «X4_LIN» and
«X10_LIN»). At $10\varepsilon$ the shear stress in the lobes reaches
«X10_SHEAR»~MPa. The calculation is static and at 0~K, so it does not include
thermally activated nucleation of dislocations. Nor does the cell contain
pre-existing dislocations, whose motion sets the yield stress of the alloy,
so the absence of defects does not mean that the alloy near such an
inclusion deforms only elastically.
